\section{7. Discussion}\label{discussion}

FLEX-DBS provides flexible, bilateral, constant-current DBS with real-time wireless control for experiments in freely moving rodents. Its measured power consumption permits a projected three days of continuous stimulation per pair of cells, with longer studies possible through a brief battery exchange. Removing the tether while retaining the ability to set, change and verify parameters at the cage allows paradigms our laboratory has used with tethered hardware \citep{creed2012egr,creed2011,aleksandrova2013,creed2013} to move into home-cage, social and extended-duration assays in which a tether is a significant confound.

The in-vivo load measurements reinforce that DBS delivery is a systems problem. Electrode impedance and stimulator compliance jointly determine whether the commanded constant current is actually delivered, and impedance differs between electrodes and varies over the weeks after implantation. A stimulator specified only by its programmable current range has not been fully specified. For the implanted platinum electrodes used here, with pulse-derived effective loads of 32--80 kΩ and weekly means of 43--54 kΩ, the ±5 V rails of FLEX-DBS support roughly 56--131 µA, well below the 600 µA programmable range and centered on the 100 µA protocols used in prior work. We therefore recommend reporting the compliance limit alongside the measured load distribution (Figure 4B) for any rodent stimulator.

Two practical consequences follow. First, electrode impedance is the most direct lever on usable amplitude, and future protocols will establish a method for fabricating lower-impedance electrodes. Second, amplitude should be titrated in each animal against its behavioral effect, as in our protocols, rather than taken from the application's nominal value. A configurable stimulator lets the experimenter confirm that a chosen amplitude is both deliverable into the implanted load and behaviorally effective, and adjust it if either condition fails.

Flexible multiday systems and fully implanted long-duration systems occupy different optimization spaces, and the measured current budget shows why. Raising the amplitude setting from 60 to 400 µA adds only 0.27 mA, whereas entering the stimulation state adds 1.68--1.95 mA over idle, so most of the current drawn during stimulation reflects circuitry and processing required while the protocol runs rather than delivered amplitude. The design carries that overhead deliberately, because each block enables a capability of FLEX-DBS: a boost converter and inverter provide ±5 V compliance; a continuously powered precision DAC and reference set the current reproducibly; an awake microcontroller sequences each pulse; and the radio stack permits configuration and verification at the cage. Power was reduced within these constraints, most visibly by shutting the output amplifier down between pulses, but maintenance-free operation for weeks or months was not the primary objective.

A fully implanted derivative would invert these priorities: battery-direct operation at $\sim$3 V if electrode impedance allows it, a shared current source with hardware-timed polarity switching, asymmetric or passive charge recovery, a radio used only to configure, and a microcontroller whose hardware peripherals can sequence a pulse without waking the core. Published mouse-scale systems show that this combination reaches average currents of roughly 8--22 µA \citep{paulat2026,plocksties2021,grotemeyer2024,fleischer2020}.

\textbf{Limitations.} Firmware v0.5 delivers the 130 Hz setting as 142.5 Hz and leaves a net charge of 26--35 \% of the stimulating-phase charge per pulse into 1 kΩ; this must be corrected before prolonged in-vivo stimulation. Below compliance, the finite output impedance of the Howland pumps (about 0.9 MΩ worst case with 0.1 \% resistors) lets delivered current fall by up to 3--8 \% into the loads measured here \citep{tucker2013}, so amplitude settings should be read against the measured load and the bench measurements rather than as absolute currents. The two channels share amplitude and timing, and the 80 Hz frequency floor excludes the low-frequency protocols of comparative and therapeutic interest \citep{creed2011,creed2015}. Finally, the bench measurements cover one pulse width into two loads, and the impedance data come from five electrodes in three mice.

\textbf{Future development.} The schematic, firmware, application and protocol are released so that others can reproduce or adapt the platform (Supplementary material). Several extensions are firmware-level within the present hardware: computing the stimulation period in timer ticks rather than whole milliseconds, which would deliver 130 Hz and other frequencies exactly; calibrating the phase-timing compensation and recovery amplitude against bench measurements so that the two phases carry equal charge; independent per-channel timing; a wider frequency range; a selectable recovery ratio; and over-the-air firmware update. At the hardware level, on-board sensing of delivered current or electrode voltage would report the stimulation actually delivered and track electrode impedance over the course of an experiment, supporting amplitude titration in each animal.

\section{8. Conclusions}\label{conclusions}

FLEX-DBS is a head-mounted, bilateral, constant-current stimulator with real-time smartphone control, intended for experiments in which parameters must be set, changed and verified at the cage. Its measured power consumption permits a projected three days of continuous operation per cell pair; firmware v0.5 requires improved charge balancing before prolonged in-vivo stimulation. Its ±5 V compliance rails deliver roughly 56--131 µA into implanted platinum electrodes with pulse-derived effective loads of 32--80 kΩ, so the electrode load, not the programmable range, sets usable amplitude in vivo.
